\documentclass[11pt]{article}

\usepackage{amsmath}
\usepackage{amssymb}
\usepackage{amsthm}
\usepackage[hidelinks]{hyperref}
\usepackage{booktabs}

\newtheorem{theorem}{Theorem}[section]
\newtheorem{lemma}[theorem]{Lemma}
\newtheorem{proposition}[theorem]{Proposition}
\newtheorem{corollary}[theorem]{Corollary}
\newtheorem{conjecture}[theorem]{Conjecture}

\theoremstyle{definition}
\newtheorem{question}[theorem]{Question}

\theoremstyle{remark}
\newtheorem{remark}[theorem]{Remark}

\newcommand{\Z}{\mathbb{Z}}
\newcommand{\Q}{\mathbb{Q}}
\newcommand{\F}{\mathbb{F}}

\title{Erd\H{o}s problem \#1056 as a moment problem\\ for factorials modulo $p$}
\author{Jeff Hinojosa}
\date{July 25, 2026}

\begin{document}

\maketitle

{\let\thefootnote\relax\footnotetext{With machine-assisted proof search and
verification. All computations reported here, and substantial portions of the
proof search and adversarial checking, were carried out by AI agents (Claude,
Anthropic) under the author's direction. Companion to \cite{Hino}, whose
Remark~6.1-style provenance disclosure applies here as well.}}

\begin{abstract}
Erd\H{o}s problem \#1056 asks, for each $k \ge 2$, for a prime $p$ and $k$
consecutive integer intervals each of product $\equiv 1 \pmod p$ ---
equivalently, for $k+1$ residues in $[1, p-1]$ sharing a common factorial
value modulo $p$. All known unconditional \emph{infinitude} results operate
at the pinned common values $\pm 1$, where Wilson-type identities apply,
and \cite{Hino} proves that toolkit terminates at $k = 3$. This note records an elementary
but consequential translation: for every $k$ there is an explicit $r = r(k)$
such that a lower bound on the $r$-th moment of the multiplicity function of
$n! \bmod p$, held on an infinite set of primes, implies \#1056 at $k$ with
infinitely many prime witnesses --- at a \emph{free} common value, entirely
outside the scope of the pinned-value obstruction. Under the standard
Poisson heuristic for factorials modulo $p$ the hypothesis holds for every
$k$ with an explicit margin. The first case not already known
unconditionally is $r(4) = 12$: a lower bound
$V_{12}(p) > 4^{11}(p-1)$ on infinitely many primes implies the infinitude
of pentad primes. We state the calibration table, the supporting census
($97.7\%$ of the primes $7 \le p < 20000$ carry $k \ge 4$ at some value;
only $15.2\%$ at value $1$), and the governing conjecture, and we pose the
minimal open analytic targets.
\end{abstract}

\section{The bridge}

Fix an odd prime $p$. For $c \in \F_p^\times$ let
\[
  \mu_p(c) \;=\; \#\{\, n \in [1, p-1] : n! \equiv c \pmod p \,\}
\]
be the multiplicity function of the factorial map, and for $r \ge 1$ let
\[
  V_r(p) \;=\; \sum_{c \in \F_p^\times} \mu_p(c)^r
  \;=\; \#\{\, (n_1, \dots, n_r) \in [1, p-1]^r :
    n_1! \equiv \dots \equiv n_r! \pmod p \,\}
\]
be its $r$-th moment. Since $n! \not\equiv 0$ for $n \le p-1$, we have
$\sum_c \mu_p(c) = p - 1$ exactly. Throughout, intervals are
$I_i = (q_{i-1}, q_i]$ with $1 \le q_0 < q_1 < \dots < q_k \le p-1$: this
is the equal-factorials form of Noll and Simmons (see \cite{EP} and
Guy \cite{Gu04}, A15), the convention recorded by OEIS A060427, and it
excludes the degenerate leading block that $0! = 1$ would otherwise
supply. Write $k(p)$ for the largest such $k$; then
$k(p) + 1 = \max_c \mu_p(c)$, as is immediate: the interval
$(q_i, q_{i+1}]$ has unit product if and only if
$q_i! \equiv q_{i+1}! \pmod p$.

\begin{lemma}[Moment bridge]\label{lem:bridge}
Let $k \ge 1$ and $r \ge 2$. If $V_r(p) > k^{r-1}(p-1)$, then
$k(p) \ge k$. Consequently, if $V_r(p) > k^{r-1}(p-1)$ for infinitely many
primes $p$, then Erd\H{o}s \#1056 holds at $k$ with infinitely many prime
witnesses.
\end{lemma}

\begin{proof}
Suppose $\max_c \mu_p(c) \le k$. Then
$\mu_p(c)^r \le k^{r-1}\mu_p(c)$ for every $c$, so
$V_r(p) \le k^{r-1} \sum_c \mu_p(c) = k^{r-1}(p-1)$. The contrapositive
gives some $c$ with $\mu_p(c) \ge k+1$: that is $k+1$ distinct elements of
$[1,p-1]$ with a common factorial value, hence $k$ consecutive intervals of
unit product. The bound is uniform in $p$, so infinitely many qualifying
$p$ give infinitely many witnesses.
\end{proof}

The lemma is pigeonhole and makes no claim to depth. Its content is the
\emph{translation}: it converts \#1056 from ``produce a rare arithmetic
configuration'' into ``beat an explicit constant in an averaged count'',
restating the problem in the language of moments and multiplicative
energy, where analytic technology operates. Two features deserve emphasis.

First, the common value $c$ produced by the lemma is \emph{free}. The
limitative theorem of \cite{Hino} shows that the classical identity toolkit
cannot pass $k = 3$ \emph{at the pinned value $1$}; the configurations
supplied by Lemma~\ref{lem:bridge} carry no value constraint and are
untouched by that obstruction. The census below shows free values are where
almost all large-multiplicity configurations actually live.

Second, the bound in the lemma is tight (all mass at multiplicity exactly
$k$ attains it), so the exponent $r$ must be chosen large enough that the
\emph{expected} moment clears the threshold. That calibration is the next
section.

\section{Calibration under the Poisson model}

The standard heuristic --- conjectured explicitly by Cobeli and Zaharescu
\cite{CZ}, generalizing Stauduhar, with the random-model analogue proved
by Cobeli, V\^aj\^aitu, and Zaharescu \cite{CVZ}, and consistent with all
computations we are aware of, including the census below --- is that for
large $p$ the multiset $\{ n! \bmod p : 1 \le n \le p-1 \}$ behaves like
$p-1$ values drawn uniformly at random from $\F_p^\times$, so that the
multiplicities $\mu_p(c)$ behave like independent Poisson random
variables of mean $1$.
The $r$-th moment of a Poisson(1) variable is the $r$-th Bell number $B_r$
(Dobi\'nski's formula), so the model predicts
\[
  V_r(p) \;\sim\; B_r \,(p-1).
\]
Define
\[
  r(k) \;=\; \min\{\, r : B_r > k^{r-1} \,\},
\]
which exists for every $k$ since $B_r$ grows super-exponentially in $r$.
Under the model, the hypothesis of Lemma~\ref{lem:bridge} at $r = r(k)$
holds for all large $p$ with margin $B_{r(k)}/k^{r(k)-1}$:

\begin{center}
\begin{tabular}{rrl}
\toprule
$k$ & $r(k)$ & margin $B_{r(k)}/k^{r(k)-1}$ \\
\midrule
 2 &  3 & $1.250$ \\
 3 &  7 & $1.203$ \\
 4 & 12 & $1.005$ \\
 5 & 19 & $1.529$ \\
 6 & 26 & $1.746$ \\
 7 & 33 & $1.476$ \\
 8 & 41 & $1.769$ \\
 9 & 49 & $1.686$ \\
10 & 57 & $1.295$ \\
\bottomrule
\end{tabular}
\end{center}

(Continuing: $r(11), \dots, r(15) = 66, 75, 84, 93, 103$.) All values were
computed exactly from the Bell triangle. Note the conspicuously thin margin
at $k = 4$: the pentad rung $r = 12$ clears its threshold by only $0.5\%$,
with $B_{12} = 4{,}213{,}597$ against $4^{11} = 4{,}194{,}304$. Any proof
attempt at $k = 4$ via $r = 12$ has essentially no room; the first
comfortable rung at $k = 4$ is $r = 13$
($B_{13}/4^{12} = 1.648$), at the price of a higher moment.

\section{What is known unconditionally}

The ladder below Lemma~\ref{lem:bridge} is remarkably short, which
calibrates how early the unconditional theory of $n! \bmod p$ still is.

\begin{itemize}
  \item \textbf{$k(p) \ge 1$ for all $p \ge 5$} is classical: Wilson's
    theorem gives $(p-2)! \equiv 1 \equiv 1!$, the universal pair.
  \item \textbf{Non-injectivity on $[2, p-1]$} --- the collision question
    with the universal pair excluded, i.e.\ the nonexistence of
    ``socialist primes'' $p > 5$. Rokowska and Schinzel \cite{RS} proved
    restrictive congruence conditions in 1960; nonexistence was verified
    computationally to $10^9$ by Trudgian \cite{Tr} (who coined the term)
    and to $10^{11}$ by Andreji\'c and Tatarevi\'c \cite{AT}. A proof for
    all $p > 5$ has recently been announced by Abramov \cite{Ab26}; the
    preprint is elementary and, to our knowledge, not yet independently
    verified.
  \item \textbf{$k(p) \ge 2$ for all large $p$ is open}, to our knowledge.
    Empirically the exceptional set is tiny and appears complete very
    early: the only odd primes $p < 20000$ with $k(p) \le 1$ are
    \[
      p \in \{3,\ 5,\ 7,\ 13,\ 19,\ 37\},
    \]
    so $k(p) \ge 2$ for every prime $41 \le p < 20000$. We know of no
    proof of $k(p) \ge 2$ valid for all large $p$.
  \item \textbf{Existence for $2 \le k \le 14$}: OEIS A060427 \cite{OEIS};
    the record witness is $a(14) = 10428007$ (each record witnesses all
    smaller $k$). An exhaustive scan by the author shows
    $a(15) > 3.2 \times 10^7$.
  \item \textbf{Infinitude for $k \le 3$}: the tetrad theorem of
    Agust\'in-Aquino and Hern\'andez Santiago \cite{AAHS}, at the pinned
    value $1$; \cite{Hino} adds the pentad criterion (along the family
    $p \equiv 3 \bmod 4$, $p \mid q! - 1$, $q \equiv 1 \bmod 6$, the
    central residue $m = (p-1)/2$ joins the tetrad
    $\{1, q, p-1-q, p-2\}$ if and only if $h(-p) \equiv 3 \bmod 4$) and
    proves the pinned-value toolkit terminates there.
  \item \textbf{Distributional results}: Klurman and Munsch \cite{KM}
    study the image of $n! \bmod p$ in short intervals, and the number of
    residue classes it misses on average over $p$, in the direction of
    the random model. Missing-value results already give second-moment
    excess: by the identity $\sum_c (\mu_p(c) - 1)^+ = \#\{\text{missed
    classes}\}$, the unconditional lower bound of Banks, Luca,
    Shparlinski, and Stichtenoth \cite{BLSS} on missed classes yields
    $V_2(p) \ge (p-1) + \Omega(\log\log p / \log\log\log p)$ on an
    infinite set of primes, and \cite{KM} improves the missed-class count
    on average (and under GRH). In the opposite direction, character-sum
    technology (Garaev--Luca--Shparlinski \cite{GLS} and its successors)
    gives \emph{upper}-bound-flavored control on congruences with $n!$.
    All known moment excesses are $o(p)$: no bound
    $V_r(p) \ge c_r (p-1)$ with a constant $c_r > 1$ appears to be known
    for any $r \ge 2$.
\end{itemize}

\section{The census, and where the truth lives}

Direct computation over the $2259$ primes $7 \le p < 20000$
(\cite{Hino}, Remark on free common values, and the verification code in
the repository):

\begin{itemize}
  \item $2207$ primes ($97.7\%$) satisfy $k(p) \ge 4$ at \emph{some}
    common value;
  \item only $344$ ($15.2\%$) do so at the value $1$, the only regime
    where any current technique operates;
  \item the maximum multiplicity grows steadily: e.g.\ $p = 977$ carries
    six equal factorials at the value $725$ ($k = 5$, $p \equiv 1 \bmod
    4$, no class number involved).
\end{itemize}

The wall at $k = 3$ is thus a wall of \emph{provability}, not of scarcity.
The Poisson model predicts far more: the maximum of $\sim p$ independent
Poisson(1) multiplicities is the classical maximum-load statistic,
\[
  \max_c \mu_p(c) \;\sim\; \frac{\log p}{\log \log p},
\]
which motivates:

\begin{conjecture}\label{conj:aa}
$k(p) \to \infty$ as $p \to \infty$. Quantitatively,
$k(p) = (1+o(1)) \log p / \log\log p$ for almost all $p$.
\end{conjecture}

Conjecture~\ref{conj:aa} implies the full Erd\H{o}s \#1056 (every $k$ is
witnessed, indeed by all sufficiently large primes outside a density-zero
set). Erd\H{o}s already suggested that witnesses exist for every $k$
\cite{Gu04}; the conjecture upgrades this to almost-all $p$ with the
Poisson maximum-load rate, completing to all multiplicities the random
model for factorials discussed in \cite{KM} and \cite{CZ}.

\section{The open targets}

We close with the precise statements the translation makes available, in
increasing order of consequence.

\begin{question}\label{q:base}
Is $k(p) \ge 2$ for all sufficiently large $p$? Equivalently: does every
large prime admit three equal factorial residues? It holds for infinitely
many $p$ \cite{AAHS}, and unconditionally for every $p \equiv 3 \bmod 4$
with $h(-p) \equiv 3 \bmod 4$, where Mordell's identity puts
$\{1, (p-1)/2, p-2\}$ at value $1$ \cite{Hino}; the open content is the
all-but-finitely-many form. (The moment version: $V_3(p) > 4(p-1)$ would
suffice for the primes on which it holds; the model predicts
$V_3(p) \sim 5(p-1)$.)
\end{question}

\begin{question}\label{q:pentads}
Is $V_{12}(p) > 4^{11}(p-1)$ for infinitely many $p$? By
Lemma~\ref{lem:bridge} this implies infinitely many pentad primes --- at a
free common value, bypassing both the class number hypothesis and the
simultaneous-divisibility hypotheses of \cite{Hino}. The same conclusion
follows from $V_{13}(p) > 4^{12}(p-1)$, where the model's margin is
$1.648$ rather than $1.005$.
\end{question}

\begin{question}\label{q:energy}
Is there \emph{any} $r \ge 2$ and any constant $c_r > 1$ with
$V_r(p) \ge c_r\,(p-1)$ for infinitely many $p$? The known excesses are
additive and small: $(2^r - 2)$ from the universal pair ($p \ge 5$), and
$\Omega(\log\log p/\log\log\log p)$ at $r = 2$ from the missed-class
bound of \cite{BLSS} via $\sum_c (\mu_p(c)-1)^+ = \#\{\text{missed
classes}\}$. All of these are $o(p)$; no constant-factor bound is known
for any $r$, while the model predicts the constant $B_r$. Any such bound
would be the first, and the method would matter more than the constant.
\end{question}

\begin{remark}[Connection to the image size]
Write $N(p) = \#\{ n! \bmod p : 1 \le n \le p-1 \}$ for the image size ---
the central quantity of \cite{KM}, long conjectured (Erd\H{o}s--Graham;
see Guy \cite{Gu04} and the discussion in \cite{KM}) to satisfy
$N(p) \sim (1 - e^{-1})\,p$. By Cauchy--Schwarz,
$V_2(p) \ge (p-1)^2 / N(p)$, so any unconditional upper bound
$N(p) \le (1-\delta)(p-1)$ would resolve Question~\ref{q:energy} at
$r = 2$ with $c_2 = 1/(1-\delta)$ --- though not conversely: a single
value of multiplicity $\sim \sqrt{cp}$ makes $V_2$ large while leaving
$N(p) = (1-o(1))(p-1)$. The missing lower bounds on moments are thus
implied by, and at most as hard as, the missing upper bounds on the
image --- a second reason the moment formulation sits squarely in the
territory of \cite{KM}.
\end{remark}

We make no claim that these are easy: lower bounds on the multiplicative
energy of $\{n!\}$ modulo $p$ are exactly what current exponential-sum
technology does not provide, and the limitative theorem of \cite{Hino}
explains why the algebraic-identity route cannot substitute for them
beyond $k=3$. The point of this note is narrower: Erd\H{o}s \#1056 is,
after Lemma~\ref{lem:bridge}, a well-posed moment problem with an explicit
finite target for each $k$, a calibrated random model that clears every
target with stated margin, and a first collision statement beyond Wilson
recently claimed by Abramov \cite{Ab26}. It is now a problem stated
entirely in the vocabulary of \cite{KM} and its antecedents.

\begin{remark}[Verification status]
The Bell numbers, the $r(k)$ table, the census figures, the exceptional
set $\{3,5,7,13,19,37\}$, and the $a(15)$ bound were each computed at
least twice in independently written code. Lemma~\ref{lem:bridge} is
elementary; no computer verification is claimed or needed for it.
\end{remark}

\begin{thebibliography}{9}

\bibitem{Ab26}
V.~M.~Abramov,
\emph{Solution to the Erd\H{o}s problem on distinct residues of
factorials},
arXiv:2604.26429 (2026). Preprint.

\bibitem{AT}
V.~Andreji\'c and M.~Tatarevi\'c,
\emph{On distinct residues of factorials},
Publ. Inst. Math. (Beograd) (N.S.) \textbf{100} (2016), 101--106.
arXiv:1603.04086.

\bibitem{BLSS}
W.~D.~Banks, F.~Luca, I.~E.~Shparlinski, and H.~Stichtenoth,
\emph{On the value set of $n!$ modulo a prime},
Turkish J. Math. \textbf{29} (2005), 169--174.

\bibitem{CVZ}
C.~Cobeli, M.~V\^aj\^aitu, and A.~Zaharescu,
\emph{The sequence $n!$ modulo $p$},
J. Ramanujan Math. Soc. \textbf{15} (2000), 135--154.

\bibitem{CZ}
C.~Cobeli and A.~Zaharescu,
\emph{Factorials modulo $p$ and the average of modular mappings},
arXiv:2011.07582 (2020).

\bibitem{AAHS}
O.~A.~Agust\'in-Aquino and J.~Hern\'andez Santiago,
\emph{Infinite tetrads of congruent factorials} (Lean-verified development
proving the infinitude of tetrad primes),
2026. Repository: \url{https://github.com/octavioalberto/tetrads}.

\bibitem{EP}
T.~F.~Bloom,
\emph{Erd\H{o}s Problem \#1056},
\url{https://www.erdosproblems.com/1056}, accessed 2026-07-25.

\bibitem{GLS}
M.~Z.~Garaev, F.~Luca, and I.~E.~Shparlinski,
\emph{Character sums and congruences with $n!$},
Trans. Amer. Math. Soc. \textbf{356} (2004), 5089--5102.

\bibitem{Gu04}
R.~K.~Guy,
\emph{Unsolved Problems in Number Theory},
3rd ed., Problems A15 and F11, Springer-Verlag, New York, 2004.

\bibitem{Hino}
J.~Hinojosa,
\emph{A class number criterion for pentads of congruent factorials},
2026. Repository:
\url{https://github.com/jeffhino/erdos-1056-pentads}.

\bibitem{KM}
O.~Klurman and M.~Munsch,
\emph{Distribution of factorials modulo $p$},
J. Th\'eor. Nombres Bordeaux \textbf{29} (2017), 169--177.
arXiv:1505.01198.

\bibitem{OEIS}
OEIS Foundation Inc.,
\emph{Sequence A060427},
The On-Line Encyclopedia of Integer Sequences,
\url{https://oeis.org/A060427}.

\bibitem{RS}
B.~Rokowska and A.~Schinzel,
\emph{Sur un probl\`eme de M. Erd\H{o}s},
Elem. Math. \textbf{15} (1960), 84--85.

\bibitem{Tr}
T.~Trudgian,
\emph{There are no socialist primes less than $10^9$},
Integers \textbf{14} (2014). arXiv:1310.6403.

\end{thebibliography}

\end{document}
